
\documentclass{article}
\usepackage{colortbl}
\usepackage{makecell}
\usepackage{multirow}
\usepackage{supertabular}

\begin{document}

\newcounter{utterance}

\centering \large Interaction Transcript for game `eqbench', experiment `eq\_bench', episode 113 with reflexion{-}Qwen3.5{-}9B.
\vspace{24pt}

{ \footnotesize  \setcounter{utterance}{1}
\setlength{\tabcolsep}{0pt}
\begin{supertabular}{c@{$\;$}|p{.15\linewidth}@{}p{.15\linewidth}p{.15\linewidth}p{.15\linewidth}p{.15\linewidth}p{.15\linewidth}}
    \# & \multicolumn{2}{c}{Player} && \multicolumn{2}{c}{Game Master} \\
    \hline

    \theutterance \stepcounter{utterance}  
    & & & \multicolumn{4}{p{0.6\linewidth}}{
        \cellcolor[rgb]{0.9,0.9,0.9}{
            \makecell[{{p{\linewidth}}}]{
                \texttt{\tiny{[P1$\langle$GM]}}
                \texttt{You have attempted this same task before and did not succeed. These are the plans you wrote after those attempts.} \\
\texttt{Trial \#1: New plan: I failed because I assigned a high score to "Superior" based on Esther's defensive words, but I ignored the explicit context clue in her opening line: "Better alone at the top than surrounded by fools at the bottom." This specific phrasing directly expresses a preference for solitude over companionship, which is the definition of feeling lonely. While she is confident (Superior), the emotional weight of the dialogue is driven by her rejection of others, making "Lonely" the primary emotion. Next time, I will prioritize the explicit textual evidence regarding her desire for isolation over her defensive tone, assigning a high score to "Lonely" (e.g., 9) and a lower score to "Superior" (e.g., 5) to reflect that her superiority is a shield against her loneliness, not the core feeling itself. I will ensure at least one score is >0 and that the highest score aligns with the most direct textual evidence.} \\
\\ 
\texttt{Superior: 5} \\
\texttt{Lonely: 9} \\
\texttt{Regretful: 0} \\
\texttt{Scared: 0} \\
\\ 
\\ 
\\ 
\texttt{Your task is to predict the likely emotional responses of a character in this dialogue:} \\
\\ 
\texttt{Ben: Well, well, well, if it isn't the Ice Queen herself. I've heard tales of your cold{-}blooded poker face.} \\
\texttt{Esther: And you must be Benny the Bluff. I've heard you're quite the charmer. } \\
\texttt{Ben: Guilty as charged, sweetheart. But don't let that fool you. I've got a few tricks up my sleeve.} \\
\texttt{Esther: I'm sure you do. But remember, Ben, this isn't a magic show. It's a game of skill. } \\
\texttt{Ben: And a little bit of luck, wouldn't you agree?} \\
\texttt{Esther: Luck is for those who can't rely on their abilities.} \\
\texttt{Ben: Spoken like a true Ice Queen. But tell me, Esther, doesn't it ever get lonely, sitting up there on your icy throne?} \\
\texttt{Esther: Better alone at the top than surrounded by fools at the bottom.} \\
\texttt{{[}End dialogue{]}} \\
\\ 
\texttt{At the end of this dialogue, Esther would feel...} \\
\texttt{Superior} \\
\texttt{Lonely} \\
\texttt{Regretful} \\
\texttt{Scared} \\
\\ 
\texttt{Give each of these possible emotions a score from 0{-}10 for the relative intensity that they are likely to be feeling each.} \\
\\ 
\texttt{You must output in the following format, including headings (of course, you should give your own scores), with no additional commentary:} \\
\\ 
\texttt{Superior: <score>} \\
\texttt{Lonely: <score>} \\
\texttt{Regretful: <score>} \\
\texttt{Scared: <score>} \\
\\ 
\\ 
\texttt{{[}End of answer{]}} \\
\\ 
\texttt{Remember: zero is a valid score, meaning they are likely not feeling that emotion. You must score at least one emotion > 0.} \\
\\ 
\texttt{Your answer:} \\
            }
        }
    }
     \\ \\

    \theutterance \stepcounter{utterance}  
    & \multicolumn{4}{p{0.6\linewidth}}{
        \cellcolor[rgb]{0.561,0.737,0.561}{
            \makecell[{{p{\linewidth}}}]{
                \texttt{\tiny{[P1$\rangle$GM]}}
                \texttt{Superior: 9} \\
\texttt{Lonely: 2} \\
\texttt{Regretful: 0} \\
\texttt{Scared: 0} \\
            }
        }
    }
    & & \\ \\

    \theutterance \stepcounter{utterance}  
    & & & \multicolumn{2}{p{0.3\linewidth}}{
        \cellcolor[rgb]{0.9,0.9,0.9}{
            \makecell[{{p{\linewidth}}}]{
                \texttt{\tiny{[GM$|$GM]}}
                \texttt{Superior: 9} \\
\texttt{Lonely: 2} \\
\texttt{Regretful: 0} \\
\texttt{Scared: 0} \\
            }
        }
    }
    & & \\ \\

    \theutterance \stepcounter{utterance}  
    & & & \multicolumn{2}{p{0.3\linewidth}}{
        \cellcolor[rgb]{0.9,0.9,0.9}{
            \makecell[{{p{\linewidth}}}]{
                \texttt{\tiny{[GM$|$GM]}}
                \texttt{Superior: 8} \\
\texttt{Lonely: 3} \\
\texttt{Regretful: 0} \\
\texttt{Scared: 0} \\
            }
        }
    }
    & & \\ \\

    \theutterance \stepcounter{utterance}  
    & & & \multicolumn{2}{p{0.3\linewidth}}{
        \cellcolor[rgb]{0.9,0.9,0.9}{
            \makecell[{{p{\linewidth}}}]{
                \texttt{\tiny{[GM$|$GM]}}
                \texttt{game\_result = LOSE} \\
            }
        }
    }
    & & \\ \\

    \theutterance \stepcounter{utterance}  
    & & \multicolumn{4}{p{0.6\linewidth}}{
        \cellcolor[rgb]{0.95,0.95,0.95}{
            \makecell[{{p{\linewidth}}}]{
                \texttt{\tiny{[P1$\leftrightarrow$P]}}
                \texttt{Post{-}episode Reflexion (not sent to GM):} \\
\texttt{New plan: I failed because I misinterpreted Esther's "Better alone..." line as a sign of confidence (Superior) rather than an admission of isolation (Lonely). The phrase explicitly contrasts being alone with being surrounded by fools, directly signaling that solitude is her current state and preference, which defines the emotion of loneliness. While her tone is defensive and arrogant, the core sentiment driving her words is the rejection of others due to a lack of genuine connection. Next time, I will prioritize the explicit semantic meaning of her key statement over her defensive tone. I will assign the highest score to "Lonely" (9) to reflect the direct textual evidence of her desire for isolation, and a lower score to "Superior" (5) to acknowledge her attitude as a coping mechanism rather than the primary feeling. I will ensure the highest score aligns with the most direct emotional descriptor in the text.} \\
\\ 
\texttt{Superior: 5} \\
\texttt{Lonely: 9} \\
\texttt{Regretful: 0} \\
\texttt{Scared: 0} \\
            }
        }
    }
    & \\ \\

\end{supertabular}
}

\end{document}
